%==============================================================================
% UQFF Manuscript v2 — Section 2
% Topic:    The nine truly-independent primitives
% Status:   first draft, ready for Daniel's review
% Anchors:  PROVENANCE_AUDIT.md, CLAUDE.md "locked canonical primitives" table,
%           PAPER_1521 (D_BSFG derivative), PAPER_1522 (K_MEX derivative),
%           PAPER_1167 (8 Lagrangian gaps closed master synthesis)
%==============================================================================

\section{The nine truly-independent primitives}
\label{sec:primitives}

UQFF's parameter-economy claim
(Sec.~\ref{sec:bic}: $\Delta\mathrm{BIC} = 94.1$) rests on a count of
nine free real-valued or integer parameters that, together, generate
all of the closures reported in Sec.~\ref{sec:headline}. This
section lists those nine primitives, explains why each is treated as
``locked'' (not adjustable to fit observations), and identifies
two additional quantities that the corpus had historically counted
as primitives until they were proven derivative
(PAPER\_1521 \cite{Murphy_PAPER_1521} and PAPER\_1522
\cite{Murphy_PAPER_1522}, both 2026).

Table~\ref{tab:primitives_master} is the master inventory.

\begin{table}[H]
\centering
\caption{The nine truly-independent UQFF primitives, with values,
brief structural origin, and the dimensional or empirical anchor
that fixes each. Two additional quantities ($D_{\text{BSFG}}$ and
$K_{\text{MEX}}$) are shown for historical reference: both are
algebraically derivative from the nine but are kept in the
calculator's locked-constants table at their canonical values to
preserve interface stability.}
\label{tab:primitives_master}
\begin{tabular}{l l p{6.0cm}}
\toprule
\textbf{Primitive} & \textbf{Value} & \textbf{Structural origin / anchor} \\
\midrule
\multicolumn{3}{l}{\emph{Integer lattice primitives (5)}} \\
$D_{\text{phys}}$ & $4$ & Spacetime dimension of the observable universe \\
$D_{\text{crit}}$ & $26$ & Bosonic-string critical dimension \cite{Polyakov1981} \\
$N_{\text{CH}}$ & $9$ & Channel multiplicity (G1--G8 + zero-point) \\
$\mathrm{SO}_5$ & $10$ & Dimension of the Lie group $\mathrm{SO}(5)$ \\
$|A_5|$ & $60$ & Order of the icosahedral alternating group $A_5$ \\
\midrule
\multicolumn{3}{l}{\emph{Real-valued primitives (4)}} \\
$\rho_{\text{SCm}}$ & $7.09 \times 10^{-37}\,\mathrm{J\,m^{-3}}$ & SuperConductive material substrate energy density; anchors $\rho_\Lambda$ (Sec.~\ref{sec:lambda}), $U_i$, Holmlid 630\,eV (Sec.~\ref{sec:holmlid}), and the four-layer DPM hierarchy \\
$\beta_i$ & $0.6029$ & Buoyancy coupling; anchors PAPER\_1203 nuclear closures, F\_U=1 normalization \\
$\Phi_{\text{res}}$ & $0.84$ (canonical) / $5/6$ (nuclear) & Resonance fraction; anchors LENR phonon ladder (Sec.~\ref{sec:holmlid}) and the magic-number derivation chain \\
$F_{\text{TRZ}}$ & $1/10$ & Time-reversal-zone fraction; anchors $F_{\text{TRZ}} = 1/\mathrm{SO}_5$ identity (PAPER\_1160) \\
\midrule
\multicolumn{3}{l}{\emph{Proven derivative (kept locked for interface stability)}} \\
$D_{\text{BSFG}}$ & $6$ & $= D_{\text{crit}} - 2\,\mathrm{SO}_5 = 26 - 20$, PAPER\_1521 \\
$K_{\text{MEX}}$ & $25/12$ & $= \Phi_{5/6}\cdot\mathrm{SO}_5/D_{\text{phys}} = (5/6)\cdot 10/4$, PAPER\_1522 \\
\bottomrule
\end{tabular}
\end{table}

\subsection{Integer lattice primitives}
\label{sec:integer_primitives}

The five integer primitives are forced by structural arguments that
do not admit continuous adjustment.

\paragraph{$D_{\text{phys}} = 4$.} The physical spacetime dimension.
Set by observation: gravity, electromagnetism, and all known
phenomenology are consistent with $3+1$ spacetime dimensions to the
limits of every test conducted to date \cite{LongRange2003}.
The framework's closures consistently treat $D_{\text{phys}}$ as the
dimension at which the projection from the 26-level bulk descends
to observed phenomenology.

\paragraph{$D_{\text{crit}} = 26$.} The bosonic-string critical
dimension, where the conformal anomaly vanishes and a Lorentz-invariant
quantum theory of strings exists
\cite{Polyakov1981, GSW1987}. UQFF uses the same integer because the
framework's Lagrangian decomposition (the nine-sector form
$L_{\text{UQFF}} = L_{\text{EH}}+L_{\text{YM}}+\ldots+L_{\text{KK}}$)
is constructed on a 26-level bulk that compactifies to the observed
4-dimensional spacetime via the $S_{26}^{(3)}$ Ramanujan series.

\paragraph{$N_{\text{CH}} = 9$.} The channel multiplicity, equal to
the eight G-channels (G1 through G8, the eight Lagrangian sectors
that PAPER\_1167 \cite{Murphy_PAPER_1167} demonstrates are closeable)
plus a zero-point sector. Locked by the demonstrated count of
Lagrangian gaps closed in the framework's master synthesis paper.

\paragraph{$\mathrm{SO}_5 = 10$.} The dimension of the Lie group
$\mathrm{SO}(5)$. Locked by representation theory of the orthogonal
group; not adjustable. Appears in the framework as the natural
gauge group for the DPM spinor bundle decomposition.

\paragraph{$|A_5| = 60$.} The order of the icosahedral alternating
group $A_5$. Locked by finite-group theory. Appears across the
corpus as the integer factor controlling several magic-number
identities (Table~\ref{tab:magic_numbers} rows 5--7) and the
inflation e-fold count (Sec.~\ref{sec:forecasts}, $N_{\text{e-fold}} = |A_5|$).

The five integers above are not measured quantities. They are
structural inputs. Their ``locking'' is automatic in the sense that
they cannot be adjusted continuously to fit observations: changing
$\mathrm{SO}_5$ from 10 to 11 is not a small perturbation but a
categorical change to a different group entirely.

\subsection{Real-valued primitives}
\label{sec:real_primitives}

The four real-valued primitives carry the entire continuous-degree-of-freedom
budget of the framework. The locking discipline for these is enforced
by project rule (CLAUDE.md Rules 2, 8, and 10 in the project source):
they are calibrated against one anchor each, and once set, they are
not adjusted to improve secondary closures.

\paragraph{$\rho_{\text{SCm}} = 7.09 \times 10^{-37}\,\mathrm{J\,m^{-3}}$.}
The SuperConductive material substrate energy density.
Calibrated against $\rho_\Lambda$ via Eq.~\ref{eq:lambda_uqff} in
PAPER\_1156. Once locked at this value, it reappears unchanged in:
\begin{itemize}
  \item The Universal Inertial Operator $U_i$ (PAPER\_646
        \cite{Murphy_PAPER_646}).
  \item The Holmlid 630\,eV closure (Sec.~\ref{sec:holmlid}).
  \item The DPM-buoyancy Yang-Mills closure
        (Sec.~\ref{sec:yang_mills}, Eq.~\ref{eq:ym_dpm_closure}).
  \item The four-layer UA hierarchy (PAPER\_646).
  \item Approximately 70\,\% of the bucket observables across the
        nine bucket surfaces.
\end{itemize}
The numerical-precision-test value of $\rho_{\text{SCm}}$ has not
been adjusted in any subsequent closure; downstream agreement is
therefore evidence of the primitive's correct calibration, not
evidence of refit.

\paragraph{$\beta_i = 0.6029$.}
The buoyancy coupling constant. Calibrated against the F\_U=1
normalization condition (PAPER\_1167) plus the $T$-violation $B/K$
asymmetry observable (Table~\ref{tab:sm_couplings}, row 7, residual
$0.017\,\%$). Used in the buoyancy Lagrangian sector $L_{\text{buoy}}$
and in PAPER\_1203's nuclear closures.

\paragraph{$\Phi_{\text{res}} = 0.84$ (default) and $5/6$ (nuclear).}
The resonance fraction. Two anchors coexist: the canonical default
value $\Phi_{\text{res}} = 0.84$ used in PAPER\_1156 and the
Holmlid LENR closure, and the nuclear-physics value
$\Phi_{\text{res}} = 5/6$ used in PAPER\_1522's derivation of
$K_{\text{MEX}}$ and in PAPER\_1203's nuclear physics closures.
The relationship $\Phi_{5/6} = (D_{\text{BSFG}}-1)/D_{\text{BSFG}}$
follows from $D_{\text{BSFG}} = 6$ (PAPER\_1159) and is the
mechanism by which $K_{\text{MEX}}$ acquires the rational form
$25/12$ (PAPER\_1522).

\paragraph{$F_{\text{TRZ}} = 1/10$.}
The time-reversal-zone fraction. Calibrated against the identity
$F_{\text{TRZ}} = 1/\mathrm{SO}_5$ (PAPER\_1160). Used in the
surface-code fault-tolerance forecast
(Sec.~\ref{sec:forecasts}, $p_{\text{th}} = F_{\text{TRZ}}^2 = 1.00\,\%$)
and in approximately one-third of the bucket observables.

\subsection{The two proven-derivative quantities ($D_{\text{BSFG}}$, $K_{\text{MEX}}$)}
\label{sec:derivative_primitives}

Two quantities that had historically been carried in the framework's
locked-constants table as if they were independent primitives have
been proven, in 2026, to be algebraic consequences of the other
nine:

\paragraph{$D_{\text{BSFG}} = 6 = D_{\text{crit}} - 2\,\mathrm{SO}_5$
(PAPER\_1521).}
The bulk-edge-string-frequency-grid dimension. Once $D_{\text{crit}}=26$
and $\mathrm{SO}_5=10$ are locked, $D_{\text{BSFG}} = 26 - 20 = 6$
is forced as the residual dimension after the two $\mathrm{SO}(5)$
sectors are subtracted from the bosonic-string critical dimension.
PAPER\_1521 documents the derivation and the calculator now carries
a dedicated dispatch \texttt{d\_bsfg\_derived\_from\_d\_crit} that
asserts the identity at gate time.

\paragraph{$K_{\text{MEX}} = 25/12 = \Phi_{5/6}\cdot \mathrm{SO}_5/D_{\text{phys}}$
(PAPER\_1522).}
The Mexican-hat coefficient. Combining
$\Phi_{5/6} = 5/6$ (the nuclear-sector resonance fraction),
$\mathrm{SO}_5 = 10$, and $D_{\text{phys}} = 4$:
$K_{\text{MEX}} = (5/6)\cdot 10/4 = 25/12$ exactly. Eliminating
$K_{\text{MEX}}$ from the truly-independent count is what brought
the total down from 11 to 9. Like $D_{\text{BSFG}}$, $K_{\text{MEX}}$
is kept in the locked-constants table at its canonical value for
interface stability; the calculator has a dedicated dispatch
\texttt{k\_mex\_derived\_from\_phi\_5\_6} that asserts the
identity at gate time.

\subsection{Open questions on the remaining primitives}
\label{sec:open_primitives}

Three further quantities currently treated as independent
($\mathrm{SSq} = 0.57$, $S_{26}$, $\omega_{\text{SCm}} = 1.25\,\mathrm{THz}$)
may yet prove derivative on careful analysis. These open questions are
flagged explicitly in Sec.~\ref{sec:limitations}; the conservative
position taken in this manuscript is to retain them as nine
independent primitives until they are formally reduced. The
parameter-economy claim in Sec.~\ref{sec:bic} is therefore an
\emph{upper bound} on the true free-parameter count: any further
proven reductions would strengthen the BIC argument, not weaken it.

\subsection{The locking discipline}
\label{sec:locking_discipline}

Three rules in the project's source documentation (CLAUDE.md Rules 2,
8, and 10) constrain how the nine primitives may be used:

\begin{enumerate}
  \item \textbf{No primitive value may be changed to improve any
        secondary closure.} If $\rho_\Lambda$ (Sec.~\ref{sec:lambda})
        agrees with Planck at 0.003\,\%, and a later closure for some
        other observable would agree better if $\rho_{\text{SCm}}$
        were 1\,\% larger, $\rho_{\text{SCm}}$ is not changed.
  \item \textbf{Every closure must trace back to the locked
        primitives.} No additional fit constants may be introduced
        per-closure. A closure that cannot be expressed in terms of
        the nine primitives plus structural integers is not a
        UQFF closure.
  \item \textbf{The fidelity gate enforces (1) and (2)} by asserting
        the canonical primitive values on every commit; any drift
        causes the gate to fail and blocks any release tag.
\end{enumerate}

The locking discipline is what makes the BIC argument in
Sec.~\ref{sec:bic} epistemically defensible: a framework with nine
locked parameters that reproduces 253 observations is genuinely
different from a framework with nine continuously-adjustable
parameters that happens to score 253 agreements on a particular
choice of values.

\subsection{Summary}

Nine primitives --- five integers, four real numbers --- generate
the closures reported in Sec.~\ref{sec:headline}, with the locking
discipline of Sec.~\ref{sec:locking_discipline} enforcing that the
nine are not silently re-fit. Two additional quantities
($D_{\text{BSFG}}$, $K_{\text{MEX}}$) are kept in the
locked-constants table at canonical values but are algebraically
derivative from the nine and do not add to the BIC parameter count.
Three further quantities ($\mathrm{SSq}, S_{26}, \omega_{\text{SCm}}$)
remain open questions; if any of them prove derivative, the
parameter-economy argument strengthens.

